\begin{pastexam}[label={px:ntu-109-comm-dsa-25}]{臺大電信所 109 學年度（B）第 25 題}
（單選；每題 2 分，答錯不倒扣）

In this exam, the height of a singleton tree is defined as 0.

A 2-3-4 tree of 255 nodes has height at least

(A)~3 \qquad (B)~4 \qquad (C)~5 \qquad (D)~6 \qquad (E)~7
\end{pastexam}

\begin{pastsoln}
答案：照字面「255 個節點」是 (B)；如果 nodes 指的是存放的鍵，是 (A)。 本書採字面的讀法，
選 (B)。

\textbf{題目的歧義。} 2-3-4 樹的一個節點可以放 $1$ 到 $3$ 個鍵，「節點數」和「鍵數」不一樣。
$255 = 4^4 - 1$ 正好是高度 $3$ 的 2-3-4 樹最多能放的鍵數，出題時可能是以鍵來想；但題目寫的是
nodes。 作答時兩種讀法都寫一行比較保險。 高度依本卷規定以邊數計（單一節點高度 $0$）。

\textbf{讀法一：255 個節點。} 2-3-4 樹所有葉節點同深度，每個節點最多 $4$ 個子節點，高度 $h$ 時
第 $i$ 層（根在第 $0$ 層）最多 $4^i$ 個節點，總共最多
\[
1 + 4 + \cdots + 4^h = \frac{4^{h+1} - 1}{3}
\]
個。 $h = 3$ 時最多 $85 < 255$，放不下；$h = 4$ 時最多 $341 \ge 255$。 高度 $4$ 也做得到：
前四層用 $1, 4, 16, 64$ 個節點（都是 4-節點），第 $4$ 層要 $255 - 85 = 170$ 個，介於
$64 \times 2 = 128$ 和 $64 \times 4 = 256$ 之間，讓 $42$ 個節點各有 $3$ 個子節點、其餘 $22$ 個
各有 $2$ 個即可（$42 \times 3 + 22 \times 2 = 170$）。 所以至少 $4$，選 (B)。

\textbf{讀法二：255 個鍵。} 每個節點最多 $3$ 個鍵，高度 $h$ 最多 $3 \cdot \frac{4^{h+1} - 1}{3} =
4^{h+1} - 1$ 個鍵（就是資結書定理~11.1 的 $n \le m^h - 1$，$m = 4$，那裡的高度以節點數計，比
這裡多 $1$）。 $h = 2$ 時最多 $63$；$h = 3$ 時最多 $255$，剛好放滿：$85$ 個節點全是 4-節點。
所以至少 $3$，選 (A)。
\end{pastsoln}
